\article{$L$ 不可区分性札记（据 R. P. Langlands 讲座）}{Notes on $L$-indistinguishability (Based on a Lecture of R. P. Langlands)}{D. Shelstad}{2}{193--203}

本文简要介绍 [4] 的结果。讨论主要限于 $\SL_2$ 的内形式（inner forms）$G$，同时尽量采用便于推广的表述，参见 [7]。假设 $F$ 是特征零的域。$F$ 为局部域（local field）时，$G(F)$ 的与一个给定表示“$L$ 不可区分”（$L$-indistinguishable）的不可约容许表示（irreducible admissible representations）只有有限多个。这些表示组成一个 $L$ 包（$L$-packet）。考察它们的特征标（characters）的适当线性组合，可以揭示包内部的结构。$F$ 为整体域（global field）时，我们考察 $G(\A)$ 的某些 $L$ 包，并根据局部表示 $\pi_v$ 在各自 $L$ 包中的位置，描述表示 $\pi=\bigotimes_v\pi_v$ 在尖点形式空间（space of cusp forms）中的重数。
\footnotetext{本工作部分受美国国家科学基金会资助，项目号 MCS76-08218。AMS（MOS）主题分类（1970）：22E50、22E55。原文版权：1979 年，美国数学会。}

\origsection{1}{$\mathcal A(T)$、$\mathcal D(T)$ 与 $\mathcal E(T)$}
设 $G$ 是定义在特征零域 $F$ 上的连通约化群（connected reductive group），$T$ 是 $G$ 中同样定义在 $F$ 上的极大环面（maximal torus）。固定 $F$ 的代数闭包 $\overline F$。令
\[
\mathcal A(T)=\{g\in G(\overline F):\Ad(g^{-1})|_T\text{ 定义在 }F\text{ 上}\},
\qquad\mathcal D(T)=T(\overline F)\backslash\mathcal A(T)/G(F).
\]
若 $g\in\mathcal A(T)$，则 $\sigma\mapsto g_\sigma=\sigma(g)g^{-1}$ 是
$\mathfrak G=\Gal(\overline F/F)$ 的取值于 $T(\overline F)$ 的连续 $1$ 上闭链（continuous $1$-cocycle）。映射 $g\mapsto(\sigma\mapsto g_\sigma)$ 诱导 $\mathcal D(T)$ 到 $H^1(\mathfrak G,T(\overline F))$ 的某个子群 $\mathcal E(T)$ 的单射；后者定义如下。令 $G_{\mathrm{sc}}$ 为 $G$ 的导出群（derived group）的单连通覆盖群（simply-connected covering group），$T_{\mathrm{sc}}$ 为 $T$ 在 $G_{\mathrm{sc}}$ 中的逆像。$\mathcal E(T)$ 就是自然同态
$H^1(\mathfrak G,T_{\mathrm{sc}}(\overline F))\to H^1(\mathfrak G,T(\overline F))$ 的像。若 $H^1(\mathfrak G,G_{\mathrm{sc}}(\overline F))=1$，特别是 $F$ 为非阿基米德局部域（nonarchimedean local field）时，$\mathcal D(T)$ 与 $\mathcal E(T)$ 相同。

当 $G$ 含有使 $\mathcal D(T)$ 非平凡的环面 $T$ 时，就出现 $L$ 不可区分性。

\origsection{2}{与 $G$ 相联系的群：局部情形（Groups attached to $G$）}
现在假设 $F$ 是局部域。固定 $F$ 的一个有限 Galois 扩张 $K$，使 $T$ 在 $K$ 上分裂。在上一节的定义中，把 $\overline F$ 换成 $K$，把 $\mathfrak G$ 换成
$\mathfrak G_{K/F}=\Gal(K/F)$。应用 Tate--Nakayama 对偶性（Tate--Nakayama duality），可将 $\mathcal E(T)$ 等同于
\[
\left\{\lambda\in X_*(T_{\mathrm{sc}}):\sum_{\sigma\in\mathfrak G_{K/F}}\sigma\lambda=0\right\}
\]
模去
\[
\left\{\lambda\in X_*(T_{\mathrm{sc}}):
\lambda=\sum_{\sigma\in\mathfrak G_{K/F}}(\sigma\mu_\sigma-\mu_\sigma),\quad
\mu_\sigma\in X_*(T)\right\}
\]
所得的商，其中 $X_*(\ )$ 表示 $\Hom(\GL_1,\ )$。$X_*(T_{\mathrm{sc}})$ 上在后一个模上平凡的拟特征（quasi-character）$\kappa$，经限制给出 $\mathcal E(T)$ 上的特征。[7] 给每个三元组 $(G,T,\kappa)$ 配以一个 $F$ 上的拟分裂群（quasi-split group），在那里记作 $H$。我们只在 $G$ 是 $\SL_2$ 的内形式时进一步讨论这一构造。

\origsection{3}{与 $\SL_2$ 的内形式相联系的群：局部情形}
设 $G$ 是 $\SL_2$ 的内形式，$F$ 是局部域；直到第 14 节我们都沿用这一假设。于是只须考虑两个群：$\SL_2$ 以及 $F$ 上一个四元数代数中范数为一的元素组成的群。可取 $\operatorname{PGL}_2(\C)$ 为 ${}^LG^\circ$（记号同 [1]），取对角子群为指定的极大环面 ${}^LT^\circ$。

固定 $G$ 中定义在 $F$ 上的极大环面 $T$。上一节考虑的拟特征 $\kappa$，此时正是 $X_*(T)$ 上的 $\mathfrak G_{K/F}$ 不变拟特征。下面固定 $X_*(T)$ 与 $X^*({}^LT^\circ)=\Hom({}^LT^\circ,\C^\times)$ 之间的同构。若 $G=\SL_2$ 且 $T$ 为对角子群，就用 $X_*(T)$ 与 $X^*({}^LT^\circ)$ 间的配对所定义的映射；若 $T$ 是 $\SL_2$ 中任意环面，就选择一次对角化，并把由此在 $X_*(T)$ 上诱导的映射与刚规定的映射复合。若 $G$ 为各向异性形式（anisotropic form），仍可将 $T$ 看成 $\SL_2$ 中的环面，以同样方式进行。

用这一同构把 $\kappa$ 转移为 $X^*({}^LT^\circ)$ 上的拟特征，再用
$\Hom(X^*({}^LT^\circ),\C^\times)$ 与 ${}^LT^\circ$ 的典范同构，把 $\kappa$ 看成 ${}^LT^\circ$ 的元素。同时，将 $\mathfrak G_{K/F}$ 在 $X_*(T)$ 上的作用转移到 $X^*({}^LT^\circ)$ 和 ${}^LT^\circ$，并用 $\sigma_T$ 表示 $\sigma\in\mathfrak G_{K/F}$ 的新作用。

分以下两种情形讨论。若 $T$ 分裂，则 $\mathfrak G_{K/F}$ 平凡作用，$\kappa$ 是 ${}^LT^\circ$ 的任意元素。若 $T$ 各向异性，设它由 $F$ 的二次扩张 $E$ 定义。把 $K$ 固定为一个足够大的有限 Galois 扩域，使它包含 $E$；$\mathfrak G_{K/F}$ 经 $\mathfrak G_{E/F}$ 作用于 $T$。令 $\sigma^\circ$ 为 $\mathfrak G_{E/F}$ 的非平凡元素，$\alpha^\vee$ 为 $T$ 在 $G$ 中的一条余根（coroot）。则
$\sigma^\circ\alpha^\vee=-\alpha^\vee$，故 $(\kappa(\alpha^\vee))^2=1$。由于 $\alpha^\vee$ 生成 $X_*(T)$，$\kappa$ 只有两种可能。非平凡的 $\kappa$ 定义 ${}^LT^\circ$ 中的元素
$\begin{pmatrix}-1&0\\0&1\end{pmatrix}_*$；这里及全文中，
$\begin{pmatrix}a&b\\c&d\end{pmatrix}_*$ 表示矩阵
$\begin{pmatrix}a&b\\c&d\end{pmatrix}\in\GL_2(\C)$ 在 $\operatorname{PGL}_2(\C)$ 中的像。
$\sigma_T$ 在 ${}^LT^\circ$ 上的作用如下：若 $\sigma$ 在
$\mathfrak G_{K/F}\to\mathfrak G_{E/F}$ 下映到平凡元，则 $\sigma_T$ 平凡作用；若映到 $\sigma^\circ$，则作用为
\[
\begin{pmatrix}a&0\\0&b\end{pmatrix}_*\longmapsto
\begin{pmatrix}b&0\\0&a\end{pmatrix}_*.
\]

作为 ${}^LT^\circ$ 的元素，$\kappa$ 对每个 $\sigma_T$ 都不变。定义 ${}^LH^\circ$ 为 $\kappa$ 在 ${}^LG^\circ$ 中的中心化子的单位连通分支。无论 $T$ 如何，$\kappa$ 平凡时 ${}^LH^\circ={}^LG^\circ$，$\kappa$ 非平凡时 ${}^LH^\circ={}^LT^\circ$。设 $\sigma\in\mathfrak G_{K/F}$。由于 ${}^LH^\circ$ 及 ${}^LT^\circ$ 都在 $\sigma_T$ 下不变，可以给 $\sigma_T$ 乘以 ${}^LH^\circ$ 的某个内自同构，得到稳定 ${}^LT^\circ$ 且固定 ${}^LT^\circ$ 在 ${}^LH^\circ$ 中的每条根（若有）的自同构 $\sigma_H$。族
$\{\sigma_H:\sigma\in\mathfrak G_{K/F}\}$ 定义半直积（semidirect product）
${}^LH={}^LH^\circ\rtimes W_{K/F}$，这里 $K/F$ 的 Weil 群 $W_{K/F}$ 经 $\mathfrak G_{K/F}$ 作用。取对偶，得到 $F$ 上的拟分裂群 $H$。具体地：

\begin{originaltheorem}{命题（Proposition）}
\begin{enumerate}
\item[(a)] 若 $\kappa$ 平凡（无论 $T$ 如何），则
${}^LH={}^LG={}^LG^\circ\times W_{K/F}$，$H=\SL_2$。
\item[(b)] 若 $T$ 分裂且 $\kappa$ 非平凡，则
${}^LH={}^LT^\circ\times W_{K/F}$，$H=T$。
\item[(c)] 若 $T$ 各向异性且 $\kappa$ 非平凡，则
${}^LH={}^LT^\circ\rtimes W_{K/F}$，其中若 $w\in W_{K/F}$ 在
$W_{K/F}\to\mathfrak G_{K/F}\to\mathfrak G_{E/F}$ 下映到 $1$，就在 ${}^LT^\circ$ 上平凡作用；若映到 $\sigma^\circ$，就以
$\diag(a,b)_*\mapsto\diag(b,a)_*$ 作用（$E,\sigma^\circ$ 同前）；并且 $H=T$。
\end{enumerate}
\end{originaltheorem}

为表示 $H$ 由 $(T,\kappa)$ 定义，记作 $H=H(T,\kappa)$。注意在定义
$X_*(T)$ 与 $X^*({}^LT^\circ)$ 的同构时所选的对角化，不影响 $H(T,\kappa)$。

\origsection{4}{$L$ 群的嵌入（Embedding $L$-groups）}
为后文所用，先固定 ${}^LH$ 到 ${}^LG$ 的嵌入 $\iota_H$。参照上述命题，在 (a) 中令 $\iota_H$ 为恒等映射，在 (b) 中为包含映射，在 (c) 中按如下方式延拓 ${}^LT^\circ\subset{}^LG^\circ$：若 $w\in W_{K/F}$ 在
$W_{K/F}\to\mathfrak G_{K/F}\to\mathfrak G_{E/F}$ 下映到 $1$，则
\[
\begin{pmatrix}a&0\\0&b\end{pmatrix}_*\times w
\longmapsto\begin{pmatrix}a&0\\0&b\end{pmatrix}_*\times w;
\]
若 $w$ 映到 $\sigma^\circ$，则
\[
\begin{pmatrix}a&0\\0&b\end{pmatrix}_*\times w
\longmapsto\begin{pmatrix}a&0\\0&b\end{pmatrix}_*
\begin{pmatrix}0&1\\1&0\end{pmatrix}_*\times w.
\]
如前所述，对任意 $F$ 上连通约化群 $G$ 都有这些群“$H$”的类似物。一般而言，${}^LH$ 未必嵌入 ${}^LG$；[7] 说明如何处理这一点（见引理 1 及最后数段）。

\origsection{5}{轨道积分：正规化（Orbital integrals: normalization）}
下面把某些轨道积分从 $G$ 转移到 $H$。先规定转移所用的正规化。固定一对 $(T,\kappa)$。分别在 $T(F)$ 和 $G(F)$ 上选 Haar 测度 $dt,dg$。若
$f\in C_c^\infty(G(F))$，$\delta\in\mathcal D(T)$，$\gamma\in T(F)$ 在 $G$ 中正则（regular），则令
\[
\Phi^\delta(\gamma,f)=\int_{h^{-1}T(F)h\backslash G(F)}
f(g^{-1}h^{-1}\gamma hg)\frac{dg}{(dt)^h},
\]
其中 $h\in\mathcal A(T)$ 代表 $\delta$，$(dt)^h$ 是由 $dt$ 经 $\Ad h$ 得到的 $h^{-1}T(F)h$ 上的测度。\footnote{\textbf{译者注（本次校订）：}按通常约定 $\Ad(h)(t)=hth^{-1}$，这里应沿 $\Ad(h^{-1})$ 推送测度，即由 $t\mapsto h^{-1}th$ 将 $dt$ 推送到积分分母中的环面上。}对
$\delta\in\mathcal E(T)-\mathcal D(T)$，规定 $\Phi^\delta(\ ,f)\equiv0$。回忆 $\mathcal A(T),\mathcal D(T),\mathcal E(T)$ 的定义见第 1 节。

当 $\kappa$ 平凡时（无论 $T$ 如何），定义
\[
\Phi^{T/\kappa}(\ ,f)=\Phi^{T/1}(\ ,f)
=\varepsilon(G)\sum_{\delta\in\mathcal E(T)}\Phi^\delta(\ ,f),
\]
其中 $G=\SL_2$ 时 $\varepsilon(G)=1$，$G$ 各向异性时 $\varepsilon(G)=-1$。称 $\Phi^{T/1}(\ ,\ )$ 为“稳定”轨道积分（stable orbital integral）。

若 $T$ 分裂且 $\kappa$ 非平凡，定义
\[
\Phi^{T/\kappa}(\gamma,f)=\frac{|\gamma_1-\gamma_2|}{|\gamma_1\gamma_2|^{1/2}}
\sum_{\delta\in\mathcal E(T)}\kappa(\delta)\Phi^\delta(\gamma,f),
\]
其中 $\gamma_1,\gamma_2$ 为 $\gamma$ 的特征值。

若 $T$ 各向异性且 $\kappa$ 非平凡，则正规化依赖两个选择。先选 $F$ 上非平凡加法特征（additive character）$\psi_F$。若 $E$ 是与 $T$ 相联系的 $F$ 的二次扩张，就按 [5] 定义 $\lambda(E/F,\psi_F)$。再选 $T(F)$ 中正则元 $\gamma^\circ$。以 $\kappa'$ 表示与 $E$ 相联系的 $F^\times$ 的二次特征。定义
\[
\Phi^{T/\kappa}(\gamma,f)=\lambda(E/F,\psi_F)
\kappa'\left(\frac{\gamma_1-\gamma_2}{\gamma_1^\circ-\gamma_2^\circ}\right)
\frac{|\gamma_1-\gamma_2|}{|\gamma_1\gamma_2|^{1/2}}
\sum_{\delta\in\mathcal E(T)}\kappa(\delta)\Phi^\delta(\gamma,f).
\]
$\gamma$ 的特征值 $\gamma_1,\gamma_2$ 以及 $\gamma^\circ$ 的特征值
$\gamma_1^\circ,\gamma_2^\circ$ 的顺序，通过固定 $T$ 的一次对角化来规定。改变 $\psi_F$ 或 $\gamma^\circ$ 的选择，至多使正规化因子改变一个符号。

\origsection{6}{轨道积分的转移（Transferring orbital integrals）}
积分 $\Phi^{T/\kappa}(\ ,f)$ 可按如下意义转移为 $H=H(T,\kappa)$ 上的稳定轨道积分。
\begin{originaltheorem}{引理（Lemma）}
若 $f\in C_c^\infty(G(F))$，则存在 $f_H\in C_c^\infty(H(F))$，使
\[
\Phi^{T/1}(\gamma,f_H)=\Phi^{T/\kappa}(\gamma,f),
\]
并且：
\begin{enumerate}
\item[(i)] 若 $G$ 各向异性且 $\kappa$ 平凡（故 $H=\SL_2$），则对分裂的 $T'$，$\Phi^{T'/1}(\ ,f_H)\equiv0$；
\item[(ii)] 若 $G=\SL_2$ 且 $\kappa$ 平凡（故 $H$ 仍为 $\SL_2$），则对所有 $T'$，
$\Phi^{T'/1}(\ ,f_H)\equiv\Phi^{T'/1}(\ ,f)$，只要 Haar 测度 $dt'$ 的选择彼此一致。
\end{enumerate}
\end{originaltheorem}
[4] 分别处理各个情形，证明了这一结论。注意当 $H=T$ 时，$\Phi^{T/1}(\ ,f_H)$ 就是 $f_H$ 自身。

[10] 对一般的实群建立了将 $\Phi^{T/\kappa}(\ ,f)$ 转移到 $H$ 的一般框架。[9] 和 [10] 对任意实群获得上述类型结果的工作有所推进。

\origsection{7}{稳定分布及一个映射（Stable distributions and a map）}
在 $G(F)$ 上的分布空间中，取稳定轨道积分的线性张成，再对逐点收敛拓扑取闭包，所得空间称为稳定分布（stable distributions）空间；这里的稳定轨道积分就是分布
$f\mapsto\Phi^{T/1}(\gamma,f)$，其中 $\gamma$ 为 $G(F)$ 的正则半单元，$T$ 为包含 $\gamma$ 的环面。

把 $\Phi^{T/\kappa}(\ ,f)$ 转移到 $H$，建立了
$C_c^\infty(G(F))$ 与 $C_c^\infty(H(F))$ 间的对应 $(f,f_H)$。将这一对应对偶化，便得到从 $H(F)$ 上的稳定分布到 $G(F)$ 上的不变分布（invariant distributions）的映射；即，若 $\Theta_H$ 为 $H(F)$ 上稳定分布，则可由
\[
\Theta(f)=\Theta_H(f_H),\qquad f\in C_c^\infty(G(F)),
\]
定义 $G(F)$ 上关于共轭不变的分布 $\Theta$。映射 $\Theta_H\mapsto\Theta$ 是下面研究 $L$ 不可区分性的关键。

用 $\chi_\pi$ 表示 $G(F)$ 的一个不可约容许表示的无穷小等价类（infinitesimal equivalence class）的特征标；用第 5 节同样的 Haar 测度正规化，把 $\chi_\pi$ 看成 $G(F)$ 的正则半单元上的函数。特征标分布 $\chi_\pi$ 的稳定性可以用以下简单准则判断。当 $G=\SL_2$ 时令 $\widetilde G=\GL_2$；当 $G$ 各向异性时令 $\widetilde G$ 为相应四元数代数的全体可逆元群。则特征标的线性组合 $\chi$ 稳定，当且仅当 $\chi$ 在 $\widetilde G(F)$ 下不变；用更便于推广的方式表述（参见 [9]），当且仅当对 $T\subset G$，$\chi$ 在 $\mathcal A(T)$ 下不变。

\origsection{8}{局部 $L$ 包与稳定特征标（Local $L$-packets and some stable characters）}
假设已对 $\widetilde G$ 证明 Langlands 对应（Langlands correspondence）；事实上，在所考察的各情形中，已证明的结果足够使用。在以下定义中也允许 $G$ 是环面。用 $\Phi(G)$ 表示 $W'_F$ 到 ${}^LG^\circ\rtimes W'_F$ 的容许同态（admissible homomorphisms）的等价类集合（记号与定义同 [1]）。每个 $\{\varphi\}\in\Phi(G)$ 对应一个由 $G(F)$ 的不可约容许表示组成的有限集 $\Pi_{\{\varphi\}}$，称为 $L$ 包。若两个不可约容许表示属于同一个 $L$ 包，就称它们 $L$ 不可区分。$G$ 是 $\SL_2$ 的内形式时，还有一个更简单的定义：$\pi_1,\pi_2$ 为 $L$ 不可区分，当且仅当存在 $g\in\widetilde G(F)$，使 $\pi_2$ 等价于 $\pi_1\circ\Ad g$。令
$\chi_{\{\varphi\}}=\sum_{\pi\in\Pi_{\{\varphi\}}}\chi_\pi$。显然：
\begin{originaltheorem}{命题（Proposition）}
$\chi_{\{\varphi\}}$ 稳定。
\end{originaltheorem}

因为 Langlands 对应已对任意连通约化实群证明 [6]，在这一情形下也可定义特征标之和 $\chi_{\{\varphi\}}$。一般说来它未必稳定；但若 $\Pi_{\{\varphi\}}$ 的成员为缓增表示（tempered，参见 [3]），则它稳定 [9]。

从第 12 节起，不必再在记号上区分容许同态
$\varphi:W'_F\to{}^LG^\circ\rtimes W'_F$ 及其等价类；我们都记为 $\varphi$，并写 $\Pi_\varphi,\chi_\varphi$ 等。

\origsection{9}{特征标恒等式：导言（Character identities: introduction）}
回到从 $H(F)$ 上稳定分布到 $G(F)$ 上不变分布的映射。若
$\{\varphi_H\}\in\Phi(H)$，则如前所述，$\chi_{\{\varphi_H\}}$ 稳定。它在 $G(F)$ 上不变分布中的像，由 $G(F)$ 的正则半单元上的某个函数 $\chi$ 表示。这个函数可用 Weyl 积分公式（Weyl integration formula）算出。记作 $\chi_{\{\varphi_H\}}\approx\chi$。

回忆第 4 节的嵌入 $\iota_H:{}^LH\hookrightarrow{}^LG$。定义
${}^LG,{}^LH,\iota_H$ 时也可以用 $W'_F$ 代替 $W_{K/F}$。作此改动后，设 $\varphi_H$ 是 $W'_F$ 到 ${}^LH$ 的容许同态，则
$\varphi=\iota_H\circ\varphi_H$ 把 $W'_F$ 映到 ${}^LG$。设 $\varphi$ 容许；回忆局部的容许性要求 $\varphi$ 的像只落在那些“与 $G$ 相关”（relevant to $G$）的 ${}^LG$ 的抛物子群（parabolic subgroups）中（参见 [1]）。这时称 $\{\varphi\}\in\Phi(G)$ 通过
$\{\varphi_H\}\in\Phi(H)$ 分解。

设 $\{\varphi\}$ 通过 $\{\varphi_H\}$ 分解。于是，自然应在 $\Pi_{\{\varphi\}}$ 中各表示的特征标的线性组合中寻找 $\chi\approx\chi_{\{\varphi_H\}}$ 的表达式。

\origsection{10}{$S^\circ\backslash S$}
下面引入描述包结构所需的群。为此，先选取等价类 $\{\varphi\}$ 中的同态
$\varphi:W'_F\to{}^LG$ 作为代表。排除 $\operatorname{sp}(2)$ [1] 及 $G(F)$ 的相应特殊表示（special representation），只考虑容许同态
$\varphi:W_{K/F}\to{}^LG$；其中 $K,{}^LG$（以及 ${}^LH,\iota_H$）同前。

定义 $S_\varphi$ 为 $\varphi$ 的像在 ${}^LG^\circ$ 中的中心化子；
$S_\varphi^\circ$ 为 $S_\varphi$ 的单位连通分支。若 $\varphi'$ 等价于 $\varphi$，则存在 $g\in{}^LG^\circ$，使
$S_{\varphi'}=gS_\varphi g^{-1}$ 且
$S_{\varphi'}^\circ=gS_\varphi^\circ g^{-1}$。

设 $\varphi=\iota_H\circ\varphi_H$，其中 $\varphi_H$ 是 $W_{K/F}$ 到 ${}^LH$ 的容许同态，$H=H(T,\kappa)$。把 $\kappa$ 看成 ${}^LT^\circ$ 的元素。回忆 ${}^LH^\circ$ 是 $\kappa$ 在 ${}^LG^\circ$ 中的中心化子的单位连通分支。按第 3 节定义，$\sigma_H$ 固定 $\kappa$，$\sigma\in\mathfrak G_{K/F}$。所以 $\kappa$ 属于 ${}^LH$ 在 ${}^LG$ 中的像的中心，从而中心化 $\varphi_H(W_{K/F})$ 在 ${}^LG$ 中的像，也就是中心化 $\varphi(W_{K/F})$。于是 $\kappa\in S_\varphi$。定义 $s_\varphi(\kappa)$，或在 $\varphi$ 已明时简记为 $s(\kappa)$，为 $\kappa$ 在
$S_\varphi^\circ\backslash S_\varphi$ 中的陪集。回忆这个商在 [3] 中已出现过。

设 $\varphi'$ 等价于 $\varphi$，且 $\varphi'=\iota_H\circ\varphi'_H$，这里
$\varphi'_H:W_{K/F}\to{}^LH$，$H=H(T',\kappa')$。于是
$\kappa'\in S_{\varphi'}$。写 $\varphi'=\Ad g\circ\varphi$，$g\in{}^LG^\circ$，则
$g^{-1}\kappa'g\in S_\varphi$。定义 $s_\varphi(\kappa')$ 为这个元素在
$S_\varphi^\circ\backslash S_\varphi$ 中的陪集。下面将证明这一商群是交换群，因而定义与 $g$ 的选择无关。

继续使用同样的 $\varphi,\varphi'$。共轭 $\Ad g$ 诱导
$S_\varphi^\circ\backslash S_\varphi$ 与
$S_{\varphi'}^\circ\backslash S_{\varphi'}$ 的同构，把
$s_\varphi(\kappa)$ 映到 $s_{\varphi'}(\kappa)$。由于两群均交换，这个同构不依赖于等价关系 $\varphi'=\Ad g\circ\varphi$ 中 $g$ 的选择。因此可把该商群及其元素 $s_\varphi(\kappa)$ 看成与 $\{\varphi\}$ 相联系的对象。

\origsection{11}{计算（Calculations）}
显式计算 $S_\varphi^\circ\backslash S_\varphi$ 及元素
$s(\kappa)=s_\varphi(\kappa)$。每个等价类只需取一个同态
$\varphi:W_{K/F}\to{}^LG$。若写
$\varphi(w)=\varphi_1(w)\times w$，$w\in W_{K/F}$，则同态
$\varphi_1:W_{K/F}\to{}^LG^\circ=\operatorname{PGL}_2(\C)$ 可提升为
$W_{K/F}$ 的二维表示 $\widetilde\varphi_1$ [7，引理 3]。

\paragraph*{(i) 设 $\widetilde\varphi_1$ 可约。}
则它通过 $\tau:W_{K/F}\to F^\times$ 分解，由 $F^\times$ 上一对拟特征
$(\mu,\nu)$ 定义。取
\[
\widetilde\varphi_1(w)=\begin{pmatrix}\mu(\tau(w))&0\\0&\nu(\tau(w))\end{pmatrix},
\qquad w\in W_{K/F}.
\]
若 $(\mu/\nu)^2\ne1$，则 $S_\varphi=S_\varphi^\circ={}^LT^\circ$；若 $\mu=\nu$，则
$S_\varphi=S_\varphi^\circ={}^LG^\circ$。然而若 $\mu/\nu$ 的阶为二，则
\[
S_\varphi={}^LT^\circ\rtimes
\left\langle1,\begin{pmatrix}0&1\\1&0\end{pmatrix}_*\right\rangle;
\]
回忆 $(\ )_*$ 表示 $(\ )$ 在 $\operatorname{PGL}_2(\C)$ 中的像。因此
$S_\varphi^\circ\backslash S_\varphi=\Z/(2)$。

无论 $\mu/\nu$ 如何，$\varphi$ 都通过 ${}^LT$ 分解，其中 $T$ 是分裂环面。相应的 $\kappa=\diag(z,1)_*$（某个 $z\ne1$）属于
$S_\varphi^\circ$，所以 $s(\kappa)$ 平凡。注意，这个同态（以及任何同态）在 $H=H(\ ,1)$ 时都通过 $H$ 分解；$s(1)$ 也平凡。

设 $\mu/\nu$ 的阶为二。在 (ii) 中将证明，与 $\varphi$ 等价的某个同态通过 ${}^LT$ 分解，这里 $T$ 是由与 $\mu/\nu$ 相联系的 $F$ 的二次扩张 $E$ 所定义的环面，并且相应的 $s(\kappa)=s_\varphi(\kappa)$ 是
$\begin{pmatrix}0&1\\1&0\end{pmatrix}_*$ 的陪集。

\paragraph*{(ii) 设 $\widetilde\varphi_1$ 通过 $W_{E/F}$ 的表示
$\widetilde\varphi_2=\Ind(W_{E/F},E^\times,\theta)$ 分解。}
这里 $E\subset K$ 是 $F$ 的二次扩张，$\theta$ 是 $E^\times$ 上的拟特征。于是 $\varphi_1$ 通过 $\widetilde\varphi_2$ 所定义的射影表示（projective representation）$\varphi_2$ 分解。设 $\sigma^\circ$ 为
$\mathfrak G_{E/F}$ 的非平凡元，定义
$\overline\theta(x)=\theta(x^{\sigma^\circ})$，$x\in E^\times$。

将 $W_{E/F}$ 显式实现为
$\{x\times\rho:x\in E^\times,\ \rho\in\{1,\sigma^\circ\}\}$，乘法为
\[
(x\times\rho)(x'\times\rho')=x(x')^\rho a_{\rho,\rho'}\times\rho\rho',
\]
其中除 $\rho=\rho'=\sigma^\circ$ 外均有 $a_{\rho,\rho'}=1$，而
$a_{\sigma^\circ,\sigma^\circ}$ 为预先选定的元素
$\alpha\in F^\times-\operatorname{Nm}_{E/F}E^\times$。于是可假设
\[
\widetilde\varphi_2(x\times1)=
\begin{pmatrix}\theta(x)&0\\0&\overline\theta(x)\end{pmatrix},
\qquad x\in E^\times,
\]
以及
\[
\widetilde\varphi_2(1\times\sigma^\circ)=\begin{pmatrix}0&z\\z&0\end{pmatrix},
\]
其中 $z$ 为选定的 $\theta(\alpha)$ 的平方根。因此
\[
\varphi_2(x\times1)=\begin{pmatrix}\theta(x)&0\\0&\overline\theta(x)\end{pmatrix}_*,
\qquad
\varphi_2(1\times\sigma^\circ)=\begin{pmatrix}0&1\\1&0\end{pmatrix}_*.
\]
设 $T$ 是与 $E$ 相联系的环面，则 $\varphi$ 通过 ${}^LT$ 分解；事实上，第 4 节中选择嵌入 $\iota_T$ 正是为了保证这一点。回忆
$T=H(T,\kappa)$，其中作为 ${}^LT^\circ$ 的元素，
$\kappa=\diag(-1,1)_*$，$s(\kappa)$ 是它在
$S_\varphi^\circ\backslash S_\varphi$ 中的陪集。

计算 $S_\varphi$ 时，先设 $(\theta/\overline\theta)^2\ne1$，则
\[
S_\varphi=\left\{1,\begin{pmatrix}-1&0\\0&1\end{pmatrix}_*\right\},
\]
故 $S_\varphi^\circ\backslash S_\varphi=\Z/(2)$；其非平凡元已写成某个 $s(\kappa)$。

现在设 $\theta/\overline\theta$ 的阶为二。此时
$S_\varphi=\varphi_1(W_{K/F})$，即
\[
S_\varphi=\left\{1,
\begin{pmatrix}-1&0\\0&1\end{pmatrix}_*,
\begin{pmatrix}0&1\\1&0\end{pmatrix}_*,
\begin{pmatrix}0&1\\-1&0\end{pmatrix}_*\right\}.
\]
因此 $S_\varphi^\circ\backslash S_\varphi=\Z/(2)\oplus\Z/(2)$。
目前已将 $\diag(-1,1)_*$ 和 $1$ 写成 $s(\kappa)$ 的形式。
为得到其余元素，考虑与 $\varphi$ 等价、通过其他 ${}^LH$ 分解的同态
$\varphi',\varphi'':W_{K/F}\to{}^LG$。

设 $E_0$ 为由 $\theta/\overline\theta$ 定义的 $E$ 的二次扩张。则
$F^\times/\operatorname{Nm}_{E_0/F}E_0^\times=\Z/(2)\oplus\Z/(2)$，所以存在两个相异的 $F$ 的二次扩域 $E',E''$，均不同于 $E$ 且包含于 $E_0$。在 $(E')^\times$ 上选拟特征 $\theta'$，使
\[
\theta'\circ\operatorname{Nm}_{E_0/E'}=\theta\circ\operatorname{Nm}_{E_0/E},
\]
并使 $\theta'/\overline{\theta'}$ 为与 $E_0/E'$ 相联系的二次特征，其中
$\overline{\theta'}(x)=\theta'(x^{\sigma'})$，
$1\ne\sigma'\in\mathfrak G_{E'/F}$。类似地选取 $\theta''$。于是
\[
\widetilde\varphi'_2=\Ind(W_{E'/F},(E')^\times,\theta'),\qquad
\widetilde\varphi''_2=\Ind(W_{E''/F},(E'')^\times,\theta'')
\]
定义 $W_{K/F}$ 的表示 $\widetilde\varphi'_1,\widetilde\varphi''_1$，它们各自等价于
$\widetilde\varphi_1$，因而定义各自等价于 $\varphi$ 的同态 $\varphi',\varphi''$。按定义 $\varphi_1,\varphi_2$ 的同样方式定义
$\varphi'_1,\varphi'_2,\varphi''_1,\varphi''_2$，并按实现 $\varphi_2$ 的同样方式实现
$\varphi'_2,\varphi''_2$。则
$\varphi_1(W_{K/F}),\varphi'_1(W_{K/F}),\varphi''_1(W_{K/F}),S_\varphi,S_{\varphi'},S_{\varphi''}$ 全部相同，并等于
\[
\left\{1,\begin{pmatrix}-1&0\\0&1\end{pmatrix}_*,
\begin{pmatrix}0&1\\1&0\end{pmatrix}_*,
\begin{pmatrix}0&1\\-1&0\end{pmatrix}_*\right\}.
\]
与此前 $\varphi$ 的情形一样，$\varphi'$（相应地 $\varphi''$）通过
${}^L(T')$（相应地 ${}^L(T'')$）分解，$T',T''$ 分别为与 $E',E''$ 相联系的环面。设
$T'=H(T',\kappa')$，$T''=H(T'',\kappa'')$。写
$\varphi=\Ad g'\circ\varphi'=\Ad g''\circ\varphi''$，$g',g''\in{}^LG^\circ$，则
\[
s(\kappa')=s_\varphi(\kappa')=g'\begin{pmatrix}-1&0\\0&1\end{pmatrix}_*g'^{-1},
\qquad
s(\kappa'')=s_\varphi(\kappa'')=g''\begin{pmatrix}-1&0\\0&1\end{pmatrix}_*g''^{-1}.
\]
初等论证表明，共轭 $\Ad g'$、$\Ad g''$ 以及
$\Ad((g')^{-1}g'')$ 都不能固定 $\diag(-1,1)_*$，故
\[
\{s(\kappa'),s(\kappa'')\}=
\left\{\begin{pmatrix}0&1\\1&0\end{pmatrix}_*,
\begin{pmatrix}0&1\\-1&0\end{pmatrix}_*\right\}.
\]
这样，$S_\varphi^\circ\backslash S_\varphi$ 的每个元素都写成了某个 $s(\kappa)$。

还须考虑一种情形：$\theta=\overline\theta$。在 $F^\times$ 上选拟特征 $\mu$，使
$\theta=\mu\circ\operatorname{Nm}_{E/F}$。则 $\varphi$ 等价于上述 (i) 型中与
$(\mu,\zeta\mu)$ 相联系的同态，其中 $\zeta$ 是由 $E/F$ 定义的 $F^\times$ 的二次特征。现在将那个同态记作 $\varphi'$。已经知道
$S_{\varphi'}^\circ\backslash S_{\varphi'}$ 有两个元素，而
$\begin{pmatrix}0&1\\1&0\end{pmatrix}_*$ 代表非平凡陪集。接下来把这个陪集写成
$s_{\varphi'}(\kappa)$。设 $T$ 是与 $E$ 相联系的环面。则 $\varphi$ 通过 ${}^LT$ 分解，相应的
$\kappa=\diag(-1,1)_*$。若 $\varphi'=\Ad g\circ\varphi$，$g\in{}^LG^\circ$，则
\[
g\begin{pmatrix}-1&0\\0&1\end{pmatrix}_*g^{-1}
=\begin{pmatrix}0&1\\1&0\end{pmatrix}_*\pmod{{}^LT^\circ},
\]
证明完成。

\paragraph*{(iii) 设 $\widetilde\varphi_1$ 为四面体型或八面体型（tetrahedral or octahedral type）。}
则 $S_\varphi=S_\varphi^\circ=1$。这是一个直接的练习。

综上：
\begin{originaltheorem}{命题（Proposition）}
\begin{enumerate}
\item[(1)] $S_\varphi^\circ\backslash S_\varphi$ 为
$1$、$\Z/(2)$ 或 $\Z/(2)\oplus\Z/(2)$ 之一。
\item[(2)] 对其每个元素 $s$，都存在一对 $(T,\kappa)$，使 $s=s(\kappa)$。
\end{enumerate}
\end{originaltheorem}

\origsection{12}{特征标恒等式：续（Character identities: continued）}
从现在起，不再在记号上区分容许同态 $\varphi:W_{K/F}\to{}^LG$ 及其等价类，都记为 $\varphi$。如已指出的，
$S_\varphi^\circ\backslash S_\varphi$ 可看成与这个等价类相联系。

设 $\varphi$ 通过 $\varphi_H$ 分解，$H=H(T,\kappa)$，令 $s=s(\kappa)$。沿用第 8、9 节记号：
\begin{originaltheorem}{引理（Lemma）}
存在整数 $\langle s,\pi\rangle$，$\pi\in\Pi_\varphi$，使
\[
\chi_{\varphi_H}\approx\sum_{\pi\in\Pi_\varphi}\langle s,\pi\rangle\chi_\pi.
\]
\end{originaltheorem}
[4] 对各个情形分别给出证明。下面列出具体结果。

\paragraph*{(A) $G=\SL_2$。}考察上一节 (i)、(ii)、(iii) 中的 $\varphi$。

(i) 若 $(\mu/\nu)^2\ne1$ 或 $\mu=\nu$，则 $\Pi_\varphi$ 只有一个元素，记为 $\pi$；
$S_\varphi^\circ\backslash S_\varphi=1$，$\langle1,\pi\rangle=1$。
若 $\mu/\nu$ 的阶为二，则 $\Pi_\varphi$ 有两个元素，记为 $\pi_1,\pi_2$。
回忆 $S_\varphi^\circ\backslash S_\varphi=\Z/(2)$；
$\langle s,\pi_1\rangle$ 和 $\langle s,\pi_2\rangle$ 是关于 $s$ 的两个特征。这些特征依赖于第 5 节正规化轨道积分时作出的选择：某些选择使
$\langle\ ,\pi_1\rangle$ 为平凡特征，另一些选择则使 $\langle\ ,\pi_2\rangle$ 为平凡特征。

(ii) $\theta=\overline\theta$ 的情形已讨论。若 $\theta/\overline\theta$ 的阶不为二，则
$\Pi_\varphi$ 有两个元素，$S_\varphi^\circ\backslash S_\varphi=\Z/(2)$，结果同 (i)。若其阶为二，则这个商群为
$\Z/(2)\oplus\Z/(2)$，$\Pi_\varphi$ 也有四个元素。设
$\Pi_\varphi=\{\pi_i:i=1,\ldots,4\}$，则可以（并且我们就这样做）正规化轨道积分，使
$\langle\ ,\pi_i\rangle$ 正好为商群的四个特征。

(iii) 此时 $\Pi_\varphi$ 只有一个元素，记为 $\pi$；商群为 $1$，
$\langle1,\pi\rangle=1$。

\paragraph*{(B) $G$ 各向异性。}
只须考虑 (ii) 中 $\theta\ne\overline\theta$ 的 $\varphi$，以及 (iii) 中的 $\varphi$。
设 $F=\R$。由于 (iii) 也被排除，$S_\varphi^\circ\backslash S_\varphi=\Z/(2)$。
但 $\Pi_\varphi$ 只有一个元素，记为 $\pi$；$\langle s,\pi\rangle$ 是关于 $s$ 的特征，依轨道积分正规化因子的选择，可为平凡或非平凡特征。

设 $F$ 非阿基米德。则 $\Pi_\varphi$ 的元素数与 $\SL_2$ 中相应 $L$ 包相同，结果也相同，唯一例外是 $\varphi$ 为 (ii) 型且
$\theta/\overline\theta$ 的阶为二的情形。此时 $\SL_2$ 中相应的包有四个元素，而
$\Pi_\varphi$ 只有一个元素 $\pi$。必须规定
$\langle1,\pi\rangle=2$，对 $s\ne1$ 规定 $\langle s,\pi\rangle=0$；特别地，
$\langle\ ,\pi\rangle$ 不是特征。

关于推广这些恒等式，我们将在即将发表的论文中报告对实群的一些进展。若 $\kappa$ 平凡，因此 $H$ 为 $G$ 的拟分裂内形式，则只要
$\Pi_\varphi$ 的成员缓增，便已知有适当的恒等式 [9]。

\origsection{13}{局部 $L$ 包的结构（Structure of local $L$-packets）}
由上一节的结果可得：
\begin{originaltheorem}{命题（Proposition）}
若 $G$ 为拟分裂形式 $\SL_2$，则配对 $\langle\ ,\ \rangle$ 把
$\Pi_\varphi$ 非典范地等同于 $S_\varphi^\circ\backslash S_\varphi$ 的对偶群（dual group）。
\end{originaltheorem}
若 $G$ 各向异性，我们只指出函数 $\langle\ ,\pi\rangle$，$\pi\in\Pi_\varphi$ 的存在。加上适当条件后，可以期待一般拟分裂实群也有类似结论（参见 [3]）。

\origsection{14}{一个整体应用（A global application）}
现在设 $F$ 为整体域，$G=\SL_2$。局部情形的群 $H$ 也有类似物：$H$ 仍可为 $G$ 中定义在 $F$ 上的极大环面，或 $G$ 自身。

假设 $K$ 是某个足够大但有限的 $F$ 的 Galois 扩域，只考虑由形如
$\Ind(W_{K/F},W_{K/E},\theta)$ 的表示定义的同态
$\varphi:W_{K/F}\to{}^LG$，其中 $E$ 是包含于 $K$ 的 $F$ 的二次扩张，
$\theta$ 是 $E$ 的不通过 $\operatorname{Nm}_{E/F}$ 分解的 Grössen 特征（Grössencharacter）。若 $T$ 为与 $E$ 相联系的环面，则 $\varphi$ 通过 ${}^LT$ 分解（与前面相同，需要 ${}^LT$ 正确嵌入 ${}^LG$）。定义 $\Pi_\varphi$ 为所有不可约容许的
$G(\A)$ 表示 $\pi=\bigotimes_v\pi_v$ 的集合，要求对每个赋值 $v$，有
$\pi_v\in\Pi_{\varphi_v}$（容许的含义同 [1]）。$\Pi_\varphi$ 是如下意义的 $L$ 包。

若 $G(\A)$ 的两个不可约容许表示
$\pi^1=\bigotimes_v\pi_v^1$、$\pi^2=\bigotimes_v\pi_v^2$ 满足：对所有赋值 $v$，
$\pi_v^1,\pi_v^2$ 都 $L$ 不可区分，而且对几乎所有 $v$，二者等价，就称
$\pi^1,\pi^2$ 为 $L$ 不可区分。$L$ 包就是这个关系的等价类。

之所以单独讨论这类 $\Pi_\varphi$，是因为 [4] 证明了：同一包中的两个表示，在
$G$ 的尖点形式空间中可能有不同的重数。一个可能以重数一出现，另一个则不出现。其余 $L$ 包的成员确实以相同重数出现（猜想该重数为零或一）。[4] 有一个公式，给出
$\Pi_\varphi$ 中表示在尖点形式中出现的重数；下面介绍这一公式。

与局部情形相同，定义 $S_\varphi$ 为 $\varphi(W_{K/F})$ 在 ${}^LG^\circ$ 中的中心化子，
$S_\varphi^\circ$ 为其单位连通分支。如第 11 节 (ii)，商群为
$\Z/(2)$ 或 $\Z/(2)\oplus\Z/(2)$。对每个赋值 $v$，$S_\varphi$ 自然嵌入
$S_{\varphi_v}$，$S_\varphi^\circ$ 自然嵌入 $S_{\varphi_v}^\circ$，从而得到自然映射
$S_\varphi^\circ\backslash S_\varphi\to S_{\varphi_v}^\circ\backslash S_{\varphi_v}$。
记号上不区分商群中的元素及其像。

回忆局部特征 $\langle\ ,\pi_v\rangle$ 依赖于正规化轨道积分时作出的选择（第 5、12 节）。此前选了 $F_v$ 上一个非平凡加法特征 $\psi_{F_v}$，并对每个在
$F_v$ 上各向异性的环面 $T$ 选择 $T(F_v)$ 的正则元
$\gamma^\circ=\gamma_v^\circ$。现在改为选 $F\backslash\A$ 上的非平凡加法特征 $\psi$，并对每个在 $F$ 上各向异性的环面 $T$ 选取
$T(F)$ 的正则元 $\gamma^\circ$。在每个 $T$ 不分裂的赋值 $v$ 处，用
$\psi_v$ 和 $\gamma^\circ$ 指定非平凡 $\kappa$ 的 $\Phi^{T/\kappa}(\ ,\ )$。

固定 $s\in S_\varphi^\circ\backslash S_\varphi$ 及 $\pi\in\Pi_\varphi$。
对几乎所有 $v$，有 $\langle s,\pi_v\rangle=1$，因此可定义
$\langle s,\pi\rangle=\prod_v\langle s,\pi_v\rangle$。于是有 [4]：
\begin{originaltheorem}{命题（Proposition）}
\begin{enumerate}
\item[(i)] $\langle s,\pi\rangle$ 不依赖于 $\psi$ 与 $\gamma^\circ$ 的选择。
\item[(ii)] $\langle\ ,\ \rangle$ 诱导从 $\Pi_\varphi$ 到
$S_\varphi^\circ\backslash S_\varphi$ 的对偶群的典范满射。
\end{enumerate}
\end{originaltheorem}
此外还有：
\begin{originaltheorem}{定理（Theorem）}
$\pi\in\Pi_\varphi$ 在尖点形式空间中出现的重数为
\[
\frac{1}{[S_\varphi^\circ\backslash S_\varphi]}
\sum_{s\in S_\varphi^\circ\backslash S_\varphi}\langle s,\pi\rangle.
\]
\end{originaltheorem}
也就是说，$\langle\ ,\pi\rangle$ 为平凡特征的那些 $\pi$ 以重数一出现，非平凡的那些则不出现。

\origsection{15}{后记（Afterword）}
更一般地，为研究与 Shimura 簇相关的问题（参见 [8]），还可以考虑满足
\[
\operatorname{Res}_{E/F}\SL_2\subset G\subset\operatorname{Res}_{E/F}\GL_2
\]
的群 $G$ 的内形式，其中 $E$ 是 $F$ 的某个有限 Galois 扩域，并以同样方式发展局部和整体理论。类似的重数公式为 [4]
\[
\frac{d_\varphi}{[S_\varphi^\circ\backslash S_\varphi]}
\sum_{s\in S_\varphi^\circ\backslash S_\varphi}\langle s,\pi\rangle,
\]
其中 $d_\varphi$ 定义如下。若 $\varphi,\varphi':W_F\to{}^LG$ 是容许同态（参见
[1，第 16 节]），并且对每个赋值 $v$，$\varphi_v$ 与 $\varphi'_v$ 等价，就称二者处处局部等价（locally equivalent everywhere）。若 $\varphi'$ 等价于
$\omega\varphi$，其中 $\omega$ 是 $W_F$ 的取值于 ${}^LG^\circ$ 的中心的连续
$1$ 上闭链，且其在每个局部群上的限制平凡，就称二者弱整体等价（weakly globally equivalent）。$d_\varphi$ 就是 $\varphi$ 的处处局部等价类中弱整体等价类的数目。

[8] 中，$L$ 不可区分性在把某些 Shimura 簇的 zeta 函数与自守 $L$ 函数联系起来时起到作用。简言之，在 [2] 所讨论的情形中，$G$ 为全实域上某个四元数代数的全体可逆元群，没有 $L$ 不可区分性，所出现的是函数
$L(s,\pi,\rho)$（$\pi$ 为 $G(\A)$ 的自守表示，$\rho$ 为 ${}^LG$ 的某个表示，按那篇讲稿固定）。在 [8] 的一般情形中，$G$ 是全体可逆元群的子群，我们至少必须考虑某些在尖点形式（乃至自守形式）中出现不同重数的 $L$ 包
$\Pi_\varphi$。此时 $\varphi:W_F\to{}^LG$ 通过
$\iota_T:{}^LT\to{}^LG$ 分解，其中 $T$ 是 $G$ 的非分裂环面。若
$\varphi=\iota_T\circ\varphi_T$、$\pi_T\in\Pi_{\varphi_T}$ 且
$\rho_T=\rho\circ\iota_T$，则
$L(s,\pi,\rho)=L(s,\pi_T,\rho_T)$。存在自然分解
$\rho_T=\rho_T^+\oplus\rho_T^-$，其中 $\rho_T^+,\rho_T^-$ 可能可约；真正相关的是函数
$L(s,\pi_T,\rho_T^\pm)$。

\begin{thebibliography}{10}
\bibitem[1]{c32-1} A. Borel, \textit{Automorphic $L$-functions}, 本论文集，第二册，27--61 页。
\bibitem[2]{c32-2} W. Casselman, \textit{The Hasse--Weil $\zeta$-function of some moduli varieties of dimension greater than one}, 本论文集，第二册，141--163 页。
\bibitem[3]{c32-3} A. W. Knapp and G. Zuckerman, \textit{Normalizing factors, tempered representations and $L$-groups}, 本论文集，第一册，93--105 页。
\bibitem[4]{c32-4} J.-P. Labesse and R. P. Langlands, \textit{$L$-indistinguishability for $\mathrm{SL}(2)$}（待刊）。
\bibitem[5]{c32-5} R. P. Langlands, \textit{On Artin's $L$-functions}, Rice Univ. Studies, vol. 56, no. 2, pp. 23--28, 1970.
\bibitem[6]{c32-6} R. P. Langlands, \textit{On the classification of irreducible representations of real algebraic groups}（待刊）。
\bibitem[7]{c32-7} R. P. Langlands, \textit{Stable conjugacy; definitions and lemmas}（待刊）。
\bibitem[8]{c32-8} R. P. Langlands, \textit{On the zeta functions of some simple Shimura varieties}（待刊）。
\bibitem[9]{c32-9} D. Shelstad, \textit{Characters and inner forms of a quasi-split group over $\R$}, Compositio Math.（待刊）。
\bibitem[10]{c32-10} D. Shelstad, \textit{Orbital integrals and a family of groups attached to a real reductive group}, Ann. Sci. École Norm. Sup.（待刊）。
\end{thebibliography}
\noindent 哥伦比亚大学（Columbia University）。
